% !TeX root = ../papers/nc-p3.tex
\section{二重二次曲線の鉛筆と滑らかな反対核からの逆判定}
\label{sec:ribbon-graph-recognition}
\begin{theorem}[二重二次曲線上の全次数認証と同型]\label{thm:conic-ribbon-as}
二二次式完全交叉$C$の被約台が滑らかな二次曲線$P\simeq\mathbb P^1$で、一般点での重複度を2とする。
$P$へ次数2で制限される二つの可逆層$L,M$と、$N=LM$の基点自由鉛筆$W$を取ると、
反転二窓は標準Koszul AS正則となる。$L\not\simeq M$なら反対階数8で、
$(C,L,M,W)$の曲線同型が添字保持同型の必要十分条件である。この層の有理商は9次元である。
$L\simeq M$なら対称完全交叉チャートに入る。
\end{theorem}
\begin{proof}
完全交叉はCohen--Macaulayなので点に支えを持つイデアルを含まない。
冪零根$J$は一般点で平方が零だから全体で$J^2=0$、$P$上で階数1無捩れである。
$\chi(\mathcal O_C)=0$から$J\simeq\mathcal O_P(-2)$となる。
滑らかなアフィン座標の持ち上げを局所的に選ぶと、貼り合わせの差は$T_P\otimes J$の導分である。
$H^1(T_P\otimes J)=H^1(\mathcal O_P)=0$なので構造層は分裂し、
$\mathcal O_C=\mathcal O_P\oplus J$と取れる。単元の完全列から、固定した制限次数での
可逆層の選択は$H^1(J)=k$の捩れとなる。非自明な制限次数零の束は、
$1$の持ち上げを妨げるこの類が非零なので大域切断を持たない。

二束の交互切断環のKoszul性を示す。束$G$の$P$への制限$G_0$が
$\mathcal O(1),\mathcal O(2)$の直和なら、$0\to G_0(-2)\to G\to G_0\to0$から
$H^1(G)=0$で$H^0(G)\to H^0(G_0)$は全射である。従って中山の補題で$G$は大域生成される。
評価核の制限は、$H^0(G_0(-2))$の次元個の自明束と
$\deg G_0$個の$\mathcal O(-1)$の直和になる。実際、制限した評価は$G_0$の評価を経由し、
各$\mathcal O(j)$の評価核は$\mathcal O(-1)^j$である。
次の$L_i$を掛ければ再び要求した制限型となる。尾部でさらに正の次数を掛けても$H^1=0$であり、
定理\ref{thm:elliptic-pencil-all-as}の線型分解がそのまま成立する。
鉛筆の一般切断は$P$の一般点で非零で正則元となる。基点自由性で第二切断は第一切断の
零因子上で単元となるので、二元の正則性とKoszul性の降下も通る。双対化層は自明であり、
同じSerre対合で有限Frobenius性を得る。

$|L|$の四つの切断として$P$上の二次単項式の持ち上げ$x_0,x_1,x_2$と
$H^0(JL)=H^0(\mathcal O_P)$の生成元$n$を取る。像の二次式は$n^2$と
$x_0x_2-x_1^2-n\ell(x)$である。後者は$H^0(JL^2)$を$n x_i$が張るため存在する。
被約台を埋め込み、$n$が各局所自明化で冪零根を生成するため閉埋込みになり、
二式の完全交叉は同じHilbert多項式でこの像に等しい。$M$側も同じである。
従って鉛筆対合で反対階数8および両側の二次核が従う。
同じ束のときは二定義式と二鉛筆切断の持ち上げが四二次式の正則列になる。

二つの埋込み二重曲線の間の同定を復元する。一般の$a\in W$の被約制限は単純な四零点を持つ。
その零因子上の環は四つの双対数環の積である。$\ell$はこの環の単元倍を除く留数対合から得られ、
各局所socleで非零である。実際、第二切断$b$はその上で単元で、随伴により$a$のCartier因子の
双対化加群は$N$の制限である。$1/b$で自明化した留数は完全対合を与える。
片側を$H^0(JL^2)=H^0(\mathcal O_P(2))$に制限すると他側は$H^0(\mathcal O_P(4))$を経由し、
四点での非零重み付き評価により対合の階数は3である。異なる二つの被約グラフの
$(2,4)$での共通線型包の階数は高々2なので、被約同定は一意である。

被約同定を固定する。第二の候補との差は、$J$でのスカラー$c$と、
$\mathcal O_P$から$J$への導分$v\in H^0(T_P\otimes J)=k$を持つ自己同型である。
二候補の乗法汎関数を比較する際は、値束を被約制限上で同一視する。
片側だけが冪零な積は$H^0(\mathcal O_P(6))$を張り、その汎関数は上の階数3により非零である。
源だけが冪零な積から両候補はこの同じ汎関数を持ち、先だけが冪零な積から$c=1$となる。
さらに被約切断の対の和$\sum p_i\otimes q_i$で$\sum p_iq_i=0$のものを取る。
持ち上げた積は冪零となるので、片側だけが冪零な積を加えて第一グラフ上の積を零にできる。
第二グラフ上には$v\sum p_i\,dq_i$だけが残る。局所枠の変更から来る零階項は
$\sum p_i q_i$を因子に持つため消える。$\mathcal O(4)$の切断の交代対だけでも
$\sum p_i\,dq_i$は$H^0(\Omega^1_P(8))=H^0(\mathcal O_P(6))$全体を張る。
アフィン座標で$(j-i)t^{i+j-1}$ ($0\le i<j\le4$)を取ればよい。
$v\ne0$なら先ほどの非零汎関数がこの全空間を消して矛盾する。よって曲線の同定は冪零方向まで一意で、
鉛筆も隣接核から戻る。

自己同型群は$\PGL_2$の持ち上げ、導分、冪零根のスカラーからなる五次元群である
（$J=\Omega^1_P$の自然作用で被約自己同型を持ち上げられる）。
二つの束のパラメータと鉛筆の次元は$1+1+12$である。一般の鉛筆の被約次数4の射は三つ以上の
分岐値を持つので、その被約安定化群は有限である。被約恒等の作用は次数$s$の束の
貼り合わせパラメータに$\lambda\mapsto c\lambda+s v$と作用する
（導分を$J=\Omega^1_P$への$d$で規格化し、二座標の遷移$t^s$の対数微分を取る）。
異なる二つの次数2のパラメータを固定すれば$c=1,v=0$となる。従って商次元は9である。
\end{proof}
\path{computations/verify_ribbon_trace.py}では二座標の遷移に一次項を含む束を使い、
異なる二束と等しい二束で両切断階数、および七次元の微分像を有理数上で検査できる。

\begin{theorem}[滑らかな反対核とグラフによる逆判定]\label{thm:smooth-apolar-graph-test}
$w\in\operatorname{Sym}^2 E\otimes\operatorname{Sym}^2 F$の隣接切断階数が6、反対階数が8とする。
両反対核の二次式が滑らかな四次完全交叉$X,Y$を切り出すとする。
$w$が楕円鉛筆の表示を持つ必要十分条件は、同定$\phi:X\to Y$のグラフへの
$(2,2)$切断の制限を通る汎関数であることである。この条件が成立すれば鉛筆の基点自由性も自動的で、
元の反転関係列は全次数で標準Koszul AS正則である。
各原点付き同定を一つ固定して平行移動を動かすとき、判定はパラメータ曲線上の次数64の
行列式線束の切断で行える。解があれば一つの被約点である。
\end{theorem}
\begin{proof}
$X,Y$上の二次切断を$A=\mathcal O_X(2),B=\mathcal O_Y(2)$と書く。
両核を割った$w$は$H^0(A)\otimes H^0(B)$上の階数8の汎関数である。
グラフ条件を仮定し、$X$上で$L=\mathcal O_X(1), M=\phi^*\mathcal O_Y(1), N=LM$とする。
$L=M$なら両者を同定したとき四切断の積は全対称で、反対と隣接の階数が等しくなるので$L\ne M$である。
$H^0(L)\otimes H^0(M)\to H^0(N)$は全射である。一般に次数4以上の二束の積の全射性は、
定理\ref{thm:elliptic-pencil-all-as}の評価核のフィルトレーションに他方の束を掛けた$H^1$消滅から従う。
グラフ上の汎関数を$\ell\in H^0(N^2)^*$とすれば、その$H^0(N)$上の積対合の核$W$は二次元である。
もし固定因子$D>0$があれば、$W$からその因子を除いた基点自由鉛筆は次数$8-\deg D\ge2$を持つ。
鉛筆の完全列で$W\otimes H^0(N)\to H^0(N^2(-D))$の核は$H^0(\mathcal O_X(D))$である。
次元を数えると積は全射となるので、$\ell$は長さ$\deg D\le6$の$D$への制限を通る。
反対対合の階数8と矛盾する。従って$W$は基点自由で、$\ell$は鉛筆の一意なtraceのスカラー倍である。
先の全次数定理を適用すればよい。逆と同定の一意性は定理\ref{thm:elliptic-pencil-reconstruction}による。

原点付きの同定を用い、$p_1,p_2:X\times X\to X$を二射影、$m$を加法とする。
$\mathcal B=p_{2*}(p_1^*A\otimes m^*B)$は階数16の束で、高次順像は零である。
各平行移動グラフへの積の制限は全射なので、$\mathcal B^*$は自明な64次元束の部分束である。
面のRiemann--Rochより$\deg\det\mathcal B=\chi(p_1^*A\otimes m^*B)=64$である。
実際、二因子の自己交叉は零で、$p_1,m$のファイバーは一点で交わる。
十六次外積に$w$を外積した全成分の同時消滅が求める条件である。いずれも$\det\mathcal B$の切断だから、
一つの恒等的に非零の成分の高々64個の零点で残りを検査すれば足りる。
異なる平行移動の線型包は前節の合併への全射により交わらない（射影線型包の意味）。
さらにグラフ$G$の法方向の一次微分は、その上で消える切断から
$H^0(G,(A\boxtimes B)|_G)$全体へ全射である。法束は自明であり、
$A\boxtimes B$を$-2G$でひねった束の$H^1$は前節と同じ豊富性で零だからである。
平行移動方向はこの法束を生成するので、非零の$w$に対して一次微分が全て消えることはない。
従って候補パラメータの閉スキームも被約である。
\end{proof}

\begin{corollary}[滑らかな反対核を保つ特殊化]\label{cor:smooth-apolar-specialization}
滑らかな楕円鉛筆から得た射影Gramテンソルの閉包内で、隣接階数6、反対階数8、
両反対核の滑らかな完全交叉性が保たれる任意の点は、実際に基点自由楕円鉛筆から得られる。
\end{corollary}
\begin{proof}
稠密な鉛筆表示から指定点への特殊化を曲線で取り、正規化し必要なら離散付値環を有限拡大する。
一定階数の反対核と二式の横断性から、両方の曲線はその上の滑らかな固有族となる。
一般ファイバーの同定は特殊ファイバーまで同定として延びる。
実際、滑らかな曲線族の正則な全空間の間の有理写像の不定点を点のブローアップで解消すると、
例外有理曲線は滑らかな種数1の特殊ファイバーの中で点に写る。
従って元の正規な面へ降下する（グラフの射影が一点ファイバーを持つ固有双有理射となる）。
逆写像も同じである。延長したグラフへの二次切断の積は両ファイバーで全射で、核は部分束である。
その核で$w$が消える条件が特殊点にも残るので、前定理を適用できる。
\end{proof}

\begin{theorem}[一つの単純零点からの楕円被覆認識]\label{thm:simple-zero-smooth-elliptic}
定理\ref{thm:smooth-apolar-graph-test}の切断階数と滑らかさのもとで、第一窓の六式に一つでも
単純零点対（六次元接空間でJacobi階数6）があれば、グラフ条件が成立する。
従ってこの場合は所属成分をあらかじめ仮定せずに全次数のAS性と楕円データの復元が得られる。
\end{theorem}
\begin{proof}
単純零点では六次元Gram空間$H$への行・列評価の像が相補三平面であり、
二重交代化の消滅により直交するので両方とも非退化である。
従って第\ref{sec:marked-opposite-rank}節の無跡対称三行列$A,B$による表示を取れる。
三列$A_i$が従属なら反対切断の最初の四行の間に関係ができ、反対核が
指定文字を因子とする可約二次式を含んでしまう。非退化な滑らかな楕円四次曲線を含む
可約二次曲面は存在しない。$B_i$も同様に独立である。

交換子行列$D_A$の階数が1以下の場合を除く。三次元の無跡対称行列空間系の中で
可換な独立二平面は、その中の正則行列（最小多項式の次数3）の無跡中心化代数に一致する。
実際、非正則な無跡対称行列の射影軌跡は$vv^{\mathsf t}$のスカラー行列を法とした像である
（ある固有値の固有空間が二次元以上）。これは二次Veronese曲面を階数3の点から射影した
埋込みであり、直線を含まない。正則行列の中心化代数はその多項式なので主張を得る。
階数0は三次元可換系を要するので不可能である。
階数1では三次元の系内に異なる二つの可換二平面があり、その共通直線は非正則行列を与える。
これは三次元で全ての二ベクトルが分解可能であることと上の中心化代数による。
共通直線の行列を$vv^{\mathsf t}$からスカラーを引いて表すと、三次元の系の全てが$v$を固有ベクトルに持つ。
従って$I,A_i$の直交補は$vu$ ($u\perp v$)という二次式の二次元空間である。
反対核の二式を指定文字の座標が零の平面に制限すると、この直交補に入る
（第\ref{sec:marked-opposite-rank}節の十行行列の最初の四行）。
従ってその平面内の直線$v=0$全体が楕円四次曲線に入ることになり矛盾する。

残る階数2のデータは可逆交換子を持つデータの閉包に属することを示す。
独立な無跡三列が張る三次元空間を固定せずに動かす。可換二平面の族は既約三次元である。
正則無跡対称行列の射影四次元開集合からその中心化二平面を取る写像が全射で、
各ファイバーは射影直線の開部分だからである。これを含む三次元空間の選択は射影二次元であり、
順序付き基底の選択は九次元である。階数2では可換二平面がただ一つとなるので、
この階数2の三列の族は既約14次元である。空でないことは対角な独立二行列と一般の第三行列で確認できる。

$A\mapsto D_A$の階数2の値には、左からの$\mathrm{SO}_3$と右からの$\mathrm{GL}_3$
（行列の同時直交変換と列変換）が作用する。右作用は余因子行列を通じるが全射である。
したがって値の軌道は像平面の法線が非等方か等方かで尽くされ、それぞれ次元8と7である。
非等方法線が実現することは、実有理数上で対角二行列と一般の実対称第三行列を取れば従う。
よって非等方軌道上のファイバー次元は6である。等方側は既約14次元の真閉部分に入るため、
その軌道が現れる場合にもファイバー次元は高々6である。
$D_B=-D_A^{\mathsf t}$により、階数2の対$(A,B)$の族の次元は高々$8+6+6=20$である。
一方、独立三列の対を動かす30次元空間ではこの等式は九式であり、各成分の次元は21以上である。
階数2以上の開部分で比較すれば、階数2だけに含まれる成分はない。

可逆交換子の全表示は命題\ref{prop:marked-net-parameters}のパラメータで覆われる。
二符号は基底変換で移り合い、射影基底を入れた像の閉包は既約である。
実際、有理な基底空間上の平方根被覆が分裂しても、二枚の像は連結な基底変換群で飽和すれば同じになる。
その一般点は楕円鉛筆の一般点である。これは同命題の支配性と
定理\ref{thm:elliptic-pencil-reconstruction}の稠密開被覆による（一般の第一窓からはGram直線が一意）。
従ってこの証明で得た階数2の極限も、基底付き射影テンソルとして楕円表示の閉包に入る。
系\ref{cor:smooth-apolar-specialization}により結論が従う。
\end{proof}

\begin{proposition}[片側の小行列式による単純零点の検査]\label{prop:smooth-apolar-defect}
上の階数と両反対核の滑らかさを仮定し、$\beta:E^*\times F^*\to H$を六次元Gram表示とする。
$B_x=\beta(x,-)$、$Q_x=B_x^{\mathsf t}h B_x$（$h$はGram計量）と書く。
第一窓の零点スキームは$\mathbb P(E^*)$内の$B_x$の四次小行列式の零点$Z$に同型である。
その各閉点で、唯一の零点対の接空間次元は$3-\operatorname{rank}Q_x\le2$である。
従って$Z$上の一点でも$Q_x$の三次小行列式が非零なら、楕円鉛筆表示とAS正則性が従う。
$Z$が有限なら、単純点の個数を$a$、接空間次元1,2の点数を$b,c$として
$a+2b+4c\le20$である。特に台に11点以上ある有限窓はこの判定で認証される。
\end{proposition}
\begin{proof}
$x\ne0$を固定する。反対核に可約二次式がないので、$x E^*$の四次元空間は
反対切断で単射に写る。その像を$F^*$上の二次形式と見ると、全てが$\ker B_x$を根基に含む。
実際、$f\in\ker B_x$なら$w(xe,ft)=\langle\beta(x,f),\beta(e,t)\rangle=0$である。
従って$\binom{\operatorname{rank}B_x+1}{2}\ge4$で階数は3以上となる。両側で同様である。
$Z$では階数3の小行列を可逆にすると、局所的な行列の行・列消去で単位三行と
残る一列三行となる。後者の三成分が四次小行列式のイデアルを生成し、
零点対のもう一方の射影座標は核直線で決まる。これはスキーム上での主張である。

零点対$(x,y)$では、三平面$U=\operatorname{im}B_x$と$V=\beta(-,y)(E^*)$は
Gram恒等式と$\beta(x,y)=0$により直交し、$V=U^\perp$である。
Jacobi像は$U+V$なので、接空間次元は$\dim(U\cap U^\perp)=3-\operatorname{rank}Q_x$となる。
反対核には$x^2$も入らず、$Q_x=w(x^2,-)$は非零である。
従って単純点の検出には$Z$を$Q_x$の三次小行列式の非消滅集合に制限すればよい。
同値な係数検査として、$B_x$の最大小行列式の斉次イデアルの根基に
$Q_x$の三次小行列式が全て属するかどうかを調べられる。

有限の場合は六つの$(1,1)$因子の固有交叉で、全長は$\binom63=20$である。
接空間次元が$k$の点では完備局所環で独立な線型部を消去すると、
$k$変数の正則局所環の極大イデアルの平方内の$k$元による零次元完全交叉となる。
$k=1$なら長さは2以上、$k=2$なら次数0,1で$1+2$次元、次数2にも
三つの単項式に対し関係が高々二本なので一次元以上残る。これで不等式を得る。
\end{proof}
